\chapter*{Özet}
\addcontentsline{toc}{chapter}{Özet}

%% Edit below this line
Kuantum hesaplamadaki gelişmelerle birlikte, RSA ve Eliptik Eğri Kriptografisi (ECC) gibi geleneksel kriptografik sistemler potansiyel güvenlik açıklarıyla karşı karşıya kalmaktadır. Bu proje, NIST tarafından standartlaştırma için seçilen, lattice tabanlı bir Anahtar Kapsülleme Mekanizması (KEM) olan CRYSTALS-Kyber’a odaklanarak, kuantum sonrası kriptografinin teori ve uygulamasını incelemektedir.

Proje, kriptografik el sıkışma protokollerinin ve klasik ECC’nin kapsamlı bir incelemesiyle başlamakta, ardından Kyber’ın matematiksel temellerinin ayrıntılı bir analizini sunmaktadır. Bu anlayışı göstermek amacıyla, güvenli anahtar değişimi için istemci-sunucu mimarisine sahip, hem ECC hem de Kyber algoritmalarını içeren tam bir kriptografik sistem uygulanmıştır.

Uygulama, polinom aritmetiği işlemleri için modern hesaplama platformlarından yararlanan çeşitli temel optimizasyonları içermekte olup, uygulama seviyesindeki işlevsellik için Python ile arayüzlenmiştir. Performans kıyaslamaları, uygun optimizasyonlarla Kyber-512’nin, kuantum direnci sağlarken ECC’ye kıyasla daha üstün hız elde ettiğini göstermektedir. Proje ayrıca, el sıkışma sürecinin ve performans metriklerinin gerçek zamanlı görselleştirilmesi için web tabanlı bir gösterge paneli içermektedir.

Bu çalışma, kuantum sonrası kriptografinin gerçek dünya ortamlarında kullanımı için hem teorik bir analiz hem de pratik bir gösterim sunmaktadır.


%% Until here
\vfill
%% Edit after {Anahtar Kelimeler:}
\textbf{Anahtar Kelimeler:} kuantum sonrası kriptografi, CRYSTALS-Kyber, eliptik eğri kriptografisi, anahtar kapsülleme mekanizması, performans optimizasyonu.
\clearpage